El Reactor Experimental Termonuclear Internacional (International Thermonuclear Experimental Reactor, ITER) és el primer projecte que estudia la possibilitat de produir energia per fusió nuclear. De totes les reaccions de fusió possibles, la reacció entre el deuteri i el triti (dos isòtops de l'hidrogen) és la més factible amb la tecnologia actual. Aquesta fusió dona ${}^{4}_{2}\mathrm{He}$ i un neutró.

\begin{apartat}{1}
A partir de les dades, digueu quants protons i quants neutrons tenen el deuteri, el triti i el ${}^{4}_{2}\mathrm{He}$. Escriviu l'equació nuclear que correspon a aquest procés de fusió.
\end{apartat}

\begin{apartat}{1}
Calculeu l'energia que s'allibera en la reacció de fusió anterior.
\end{apartat}

\dades{%
  $1\ \mathrm{eV} = 1,602\times 10^{-19}\ \mathrm{J}$.\\
  $c = 3,00\times 10^{8}\ \mathrm{m\,s^{-1}}$.\\
  Masses (en kg):}

{\centering\renewcommand{\arraystretch}{1.4}
\begin{tabular}{|c|c|c|c|}
\hline
${}^{1}_{0}\mathrm{n}$ (neutró) & ${}^{4}_{2}\mathrm{He}$ (heli) & ${}^{2}_{1}\mathrm{H}$ (deuteri) & ${}^{3}_{1}\mathrm{H}$ (triti)\\
\hline
$1,674\,927\times 10^{-27}$ & $6,644\,657\times 10^{-27}$ & $3,343\,584\times 10^{-27}$ & $5,007\,357\times 10^{-27}$\\
\hline
\end{tabular}\par}
